% !TeX program = lualatex
\documentclass[11pt,a4paper]{article}
\usepackage[T1]{fontenc}
\usepackage[utf8]{inputenc}
\usepackage[russian]{babel}
\usepackage{amsmath,amssymb,amsfonts}
\usepackage{amsthm}                     % <- ДОБАВЛЕНО
\usepackage{geometry}
\geometry{left=1.25cm,right=1.25cm,top=1.25cm,bottom=3.1cm}
\usepackage{booktabs}
\usepackage{array}
\usepackage{hyperref}
\usepackage{url}
\usepackage{graphicx}
\usepackage{caption}
\usepackage{subcaption}
\usepackage{float}
\usepackage{siunitx}
\usepackage{bm}
\usepackage{pgfplots}
\pgfplotsset{compat=1.18}
\usepackage{xcolor}
\usepackage{titlesec}
\usepackage{fancyhdr}
\usepackage{microtype}
\usepackage{multicol}
\usepackage{fontspec}
\usepackage{tikz}
\usetikzlibrary{shapes,arrows,arrows.meta,positioning,decorations.pathmorphing,patterns,calc}
\usepackage{enumitem}

\tikzset{>=stealth}

% --------------------------------------------------------------
% Шрифты Windows (стандартные, с поддержкой кириллицы)
% --------------------------------------------------------------
\setmainfont{Times New Roman}[
Ligatures = TeX,
Script = Cyrillic,
Language = Russian
]
\setsansfont{Arial}[
Ligatures = TeX,
Script = Cyrillic,
Language = Russian
]
\setmonofont{Courier New}[
Script = Cyrillic,
Language = Russian
]
% --------------------------------------------------------------

\definecolor{skyblue}{RGB}{0,102,204}
\definecolor{darkblue}{RGB}{0,0,139}
\definecolor{gold}{RGB}{255,215,0}

\newcommand{\eps}{\varepsilon}
\newcommand{\kappac}{\kappa_c}
\newcommand{\Keff}{K_{\mathrm{eff}}}
\newcommand{\betaf}{\beta}
\newcommand{\df}{d_f}

\titleformat{\section}{\LARGE\bfseries\color{darkblue}}{\thesection}{1em}{}
\titleformat{\subsection}{\normalsize\bfseries\color{darkblue}}{\thesubsection}{1em}{}
\titleformat{\subsubsection}{\normalsize\itshape}{\thesubsubsection}{1em}{}

% Определения теорем
\newtheorem{axiom}{Аксиома}
\newtheorem{remark}{Замечание}
\newtheorem{theorem}{Теорема}[section]
\DeclareMathOperator*{\Tr}{Tr}

\fancypagestyle{firstpage}{%
	\fancyhf{}%
	\renewcommand{\headrulewidth}{0pt}%
	\renewcommand{\footrulewidth}{0pt}%
	\fancyfoot[C]{%
		\begin{minipage}[t]{\textwidth}
			\centering
			\textcolor{gold}{\rule{\textwidth}{1.6pt}}\\[5pt]
			{\small \textsuperscript{1}Yakushev Research, \url{https://yuct.org/}} \hfill
			{\small YPSDC-протокол: \url{https://ypsdc.com/}}\\[-2pt]
			\textcolor{gold}{\rule{\textwidth}{1.6pt}}\\[5pt]
			\textbf{ЯКУШЕВ ЕДИНАЯ КООРДИНАЦИОННАЯ ТЕОРИЯ 2026} \hfill \thepage\\[10pt]
		\end{minipage}%
	}%
}

\fancypagestyle{plain}{%
	\fancyhf{}%
	\renewcommand{\headrulewidth}{0pt}%
	\renewcommand{\footrulewidth}{0pt}%
	\fancyfoot[C]{%
		\begin{minipage}[t]{\textwidth}
			\centering
			\textcolor{gold}{\rule{\textwidth}{1.6pt}}\\[4pt]
			\textbf{ЯКУШЕВ ЕДИНАЯ КООРДИНАЦИОННАЯ ТЕОРИЯ 2026} \hfill \thepage\\[4pt]
		\end{minipage}%
	}%
}

\pagestyle{plain}
\setlength{\parindent}{0pt}
\setlength{\parskip}{1ex}
\raggedbottom

\hypersetup{
	colorlinks=true,
	linkcolor=blue,
	citecolor=blue,
	urlcolor=blue,
}

\newcommand{\makeheader}{%
	\noindent
	\begin{minipage}{0.17\textwidth}
		\raggedright
		{\sffamily\bfseries\color{skyblue}\fontsize{46}{46}\selectfont YUCT}%
	\end{minipage}%
	\hspace{30pt}
	\begin{minipage}{0.32\textwidth}
		\raggedright
		{\sffamily\fontsize{11}{13}\selectfont ЯКУШЕВ\\ ЕДИНАЯ\\ КООРДИНАЦИОННАЯ ТЕОРИЯ}%
	\end{minipage}%
	\hfill
	\begin{minipage}{0.38\textwidth}
		\raggedleft
		{\sffamily\bfseries Техническая заметка}\\
		{\sffamily Опубликовано апрель 2026}\\
		{\sffamily\footnotesize \scriptsize\url{https://doi.org/10.5281/zenodo.18444598}}%
	\end{minipage}
	\par\vspace{5pt}
	\noindent\textcolor{gold}{\rule{\textwidth}{1.6pt}}
	\par\vspace{0pt}
}


\begin{document}
	
		\makeheader
		\centering
		\vspace*{1cm}
		
		{\Huge\bfseries \(K_{\mathrm{eff}} > 1\) как онтологический первый принцип\par}
		\vspace{0.3cm}
		{\Large\bfseries Универсальный скейлинговый показатель \(\beta = 2/3\), электрослабая шкала\par}
		\vspace{0.1cm}
		{\Large\bfseries и квантование масс фермионов\par}
		\vspace{0.5cm}
		{\large Прозрачный синтез аксиом, эмпирических данных и феноменологии\par}
		
		\vspace{1cm}
		
		{\large Алексей В. Якушев\par}
		\vspace{0.3cm}
		{\normalsize Yakushev Research, \url{https://yuct.org/}\par}
		
		\vspace{1.0cm}
		
		{\normalsize Версия 6.3 (с гипотезой о массе нейтрино), апрель 2026\par}
		\vspace{0.3cm}
		{\normalsize DOI: \texttt{10.5281/zenodo.18444598}\par}
		
		\vfill
		
		\begin{abstract}
			\noindent
			Мы представляем окончательную логическую структуру Единой теории координации Якушева (YUCT).  Каждое утверждение классифицируется как \textbf{Аксиома} (несводимый постулат), \textbf{Теорема} (доказанная из аксиом), \textbf{Эмпирический закон} (надёжно наблюдаемый) или \textbf{Феноменологическая подгонка} (ограниченная, но ещё не выведенная).  Из определения эффективности координации \(K_{\mathrm{eff}} > 1\) и трёх структурных аксиом мы доказываем, что фактор избыточности словаря равен \(q = 2/3\).  Универсальный показатель \(\beta = 2/3\) является \textbf{эмпирическим законом}, подтверждённым более чем десятком независимых экспериментов на атомных, биологических, геофизических и космологических масштабах.  Из \(\beta = 2/3\) мы выводим алгебраические константы \(S_{\mathrm{odd}} = 1.2\) и \(S_{\mathrm{even}} = 0.8\) (замкнутая алгебраическая петля).  Используя число степеней свободы вейлевских фермионов (\(96\)), мы получаем электрослабую шкалу (\(m_W \approx 80.4\) ГэВ) с точностью до единственного множителя \(\xi_W\).  Та же алгебраическая структура даёт квант масштабирования \(q = (3/2)^{1/3}\).  Массы девяти заряженных фермионов лежат на полуцелой лестнице с точностью лучше \(1\%\) (конкретные квантовые числа \(N_f\) в настоящее время подогнаны).  Дискретная структура этой лестницы является \textbf{фальсифицируемым предсказанием}: любая будущая частица с массой вне разрешённого множества \(\{m_e q^{N/2}\}\) опровергнет теорию.  Для легчайшего нейтрино разрешённые массы составляют \(\{0.0234, 0.0268, 0.0307, 0.0352, \dots\}\) эВ.  Теория не содержит настраиваемых непрерывных параметров и делает несколько других конкретных, проверяемых предсказаний.
			
			Кроме того, мы выдвигаем вариационную гипотезу, которая выбирает конкретное квантовое число нейтрино \(N_{\nu} = -125\) (масса \(\approx 0.023\) эВ) путём максимизации полной координационной мощности фермионного сектора.  Это предсказание фальсифицируемо и, в случае подтверждения, усилит теорию.
		\end{abstract}
		
		\vspace{1cm}
		\textcopyright\ 2026 Алексей В. Якушев. Все права защищены.
	
	
	\tableofcontents
	\newpage
	
	\section{Введение: координация против информации}
	\label{sec:intro}
	
	Теория связи Шеннона ограничивает скорость, с которой информация может передаваться по каналу.  Тем не менее природа полна процессов, в которых короткий сигнал вызывает реакцию, информационное содержание которой намного превышает содержание самого сигнала: молекула феромона запускает сложное поведение, стартовый кодон инициирует синтез белка, короткая команда активирует запланированную операцию.  Дополнительная информация не передаётся в реальном времени; она уже хранится в \textbf{словаре}, распространённом во время автономной фазы.
	
	Протокол синхронной распределённой координации Якушева (YPSDC) формализует этот универсальный механизм:
	\begin{enumerate}[label=(\roman*)]
		\item \textbf{Автономная фаза:} Словарь \(\mathcal{D}\), содержащий \(M\) возможных действий (состояний, сообщений), разделяется между агентами.  Каждое действие \(A_i\) требует \(n_i\) бит для полного описания.
		\item \textbf{Онлайн-фаза:} Чтобы активировать действие \(A_k\), передаётся только его индекс \(\kappa_k\).  При равновероятных индексах длина индекса составляет \(\ell = \log_2 M\) бит.
	\end{enumerate}
	\textbf{Эффективность координации} определяется как
	\begin{equation}
		\boxed{K_{\mathrm{eff}} = \frac{H(\mathcal{A})}{H(\kappa)} = \frac{\text{Энтропия действия}}{\text{Энтропия индекса}} \ge 1},
		\label{eq:keff-def}
	\end{equation}
	где \(H\) обозначает энтропию Шеннона.  Ни один физический закон не нарушается: индекс передаётся со скоростью \(v \le c\); дополнительная информация поставляется локально из уже существующего словаря.
	
	\begin{center}
		\fbox{%
			\begin{minipage}{0.9\textwidth}
				\textbf{Онтологический принцип:} \(K_{\mathrm{eff}} > 1\) является необходимым условием для того, чтобы сложная система могла сохранять свою идентичность во времени.  Система с \(K_{\mathrm{eff}} = 1\) всегда отставала бы от своего окружения и была бы разрушена.  Следовательно, сама реальность должна быть структурирована как координационная сеть.
			\end{minipage}%
		}
	\end{center}
	
	\section{Аксиомы иерархических координационных сетей}
	
	Координационная сеть строится из вложенных кластеров.
	
	\begin{axiom}[Иерархическая масштабная инвариантность] \label{ax:A1}
		Кластер уровня \(n\) имеет линейный размер \(L_n\) с \(L_{n+1} = \lambda L_n\) (\(\lambda > 1\)).  Число агентов масштабируется как \(N(L) \propto L^{d_f}\), где \(d_f\) – фрактальная размерность.  При \(n \to \infty\) отношение \(K_{\mathrm{eff},n+1}/K_{\mathrm{eff},n}\) стремится к константе, что подразумевает степенную иерархию.
	\end{axiom}
	
	\begin{theorem}[Минимальная энтропия словаря]
		\label{thm:minimal-dict}
		Минимальный координационный словарь, способный надёжно различать бинарную альтернативу, обладает полной энтропией Шеннона ровно \(2\) бита (в единицах \(k_B\ln 2\)).  Следовательно, полный иерархический словарь удовлетворяет условию \(S_{\mathrm{odd}} + S_{\mathrm{even}} = 2\).
	\end{theorem}
	
	\begin{proof}
		Минимальному словарю достаточно отвечать на единственный вопрос «да/нет»: «Произошло ли действие \(A\)?»  Поэтому он содержит две записи, \(\mathcal{A} = \{A_0, A_1\}\), с равными априорными вероятностями, так что
		\[
		H(\mathcal{A}) = \log_2 2 = 1\ \text{бит}.
		\]
		Чтобы активировать любую из них, передаётся индекс \(\kappa \in \{0,1\}\) длины \(\ell = \log_2 2 = 1\)~бит, что даёт \(H(\mathcal{K}) = 1\)~бит.
		
		Однако приёмник также должен быть уверен, что индекс правильно идентифицирует намеченное действие; в противном случае единственный переворот бита разрушил бы координацию.  Эта уверенность требует одного дополнительного проверочного бита, хранящегося в словаре.  Следовательно, полная энтропия словаря равна
		\[
		S_{\min} = H(\mathcal{A}) + H(\text{проверка}) = 1 + 1 = 2\ \text{бита}.
		\]
		В бесконечной иерархии (раздел~\ref{sec:algebraic-loop}) эта минимальная структура реализуется суммами \(S_{\mathrm{odd}} + S_{\mathrm{even}} = 2\), что фиксирует абсолютный масштаб без каких-либо настраиваемых параметров.
	\end{proof}
	
	\begin{axiom}[Вложение в три измерения] \label{ax:A3}
		Агенты находятся в пространстве размерности \(D = 3\).  Время координации для кластера размера \(L\) составляет \(\tau(L) = L/c\) и ограничено скоростью света.
	\end{axiom}
	
	
	\section{Теорема: фактор избыточности \(q = 2/3\)}
	
	Пусть энтропия словаря на самом глубоком уровне равна \(H_1\).  Согласно аксиоме \ref{ax:A1}, энтропии последовательных уровней масштабируются геометрически: \(H_{l+1} = q H_l\) с \(0 < q < 1\).  Полная энтропия бесконечной иерархии равна
	\[
	\sum_{l=1}^{\infty} H_l = \frac{H_1}{1-q}.
	\]
	Теорема \ref{thm:minimal-dict} требует, чтобы эта сумма равнялась \(2\).  Мы дополнительно постулируем, что \(H_1 = q\) (первый уровень несёт информацию одного фактора избыточности; это простейший самосогласованный выбор).  Тогда
	\begin{equation}
		\frac{q}{1-q} = 2 \quad\Longrightarrow\quad \boxed{q = \frac{2}{3}}.
		\label{eq:q}
	\end{equation}
	Таким образом, фактор избыточности координационного словаря точно равен \(2/3\).
	
	\section{Универсальный закон ошибки: \(\varepsilon \propto K_{\mathrm{eff}}^{-\beta}\) с \(\beta = 2/3\)}
	
	Обширные эмпирические исследования (Приложение L, Ferra1.2 и сводные документы валидации) показали, что в любой координированной системе относительная ошибка \(\varepsilon\) – вероятность активации неверной записи словаря – следует степенному закону:
	\begin{equation}
		\boxed{\varepsilon = \kappa_c \, \alpha \, K_{\mathrm{eff}}^{-\beta}, \qquad \beta = \frac{2}{3} \approx 0.6667,\quad \kappa_c = \frac{1}{3}}.
		\label{eq:error-law}
	\end{equation}
	Константа \(\kappa_c = 1/3\) следует из триадной онтологии \(\text{Реальность} = \mathbf{D} + \mathbf{I} \times \mathbf{R}\) (Словарь, Информация, Резонанс), которая подразумевает минимальную координационную энтропию \(S_{\min} = k_B \ln 3\) и, следовательно, \(\kappa_c = e^{-S_{\min}/k_B} = 1/3\).
	
	Показатель \(\beta = 2/3\) \textbf{не} выводится из трёх вышеприведённых аксиом; это \textbf{эмпирически установленная универсальная константа}.  Таблица \ref{tab:beta-evidence} содержит репрезентативную выборку независимых измерений.  Полный каталог включает более десятка различных физических, биологических и социальных систем, которые все дают показатели, согласующиеся с \(2/3\) (подробный анализ см. в технических заметках и приложениях YUCT).
	
	\begin{table}[htbp]
		\centering
		\caption{Выборка экспериментальных подтверждений универсального показателя \(\beta = 2/3\).  Все приведённые значения являются наблюдаемыми скейлинговыми показателями; они согласуются с предсказанием YUCT в пределах указанных неопределённостей.}
		\label{tab:beta-evidence}
		\small
		\begin{tabular}{l c c c}
			\toprule
			\textbf{Система} & \textbf{Наблюдаемая} & \textbf{Наблюдаемый показатель} & \textbf{Ссылка} \\
			\midrule
			Энергии связей галогеноводородов & \(E \propto r^{-\beta}\) & \(0.682 \pm 0.012\) & Приложение~C \\
			Частота ошибок репликации ДНК & \(\varepsilon\) отн.\ \(K_{\mathrm{eff}}\) & \(0.65\text{--}0.70\) & Приложение~L \\
			Частота мутаций у позвоночных & \(\mu \propto K_{\mathrm{ремонта}}^{-\beta}\) & \(0.67 \pm 0.05\) & Приложение~E \\
			Метаболический скейлинг (простые организмы) & \(B \propto M^{\beta}\) & \(0.68 \pm 0.03\) & Приложение~D \\
			Аномальная диффузия в пористых средах & \(\langle r^2 \rangle \propto t^{\beta}\) & \(0.67 \pm 0.04\) & Bordoloi et al.\ (2022) \\
			Фрагментация базальта & \(N(>m) \propto m^{-\beta}\) & \(0.66 \pm 0.03\) & Hartmann (1969) \\
			Релаксация стекла вблизи \(T_g\) & \(\tau \propto (T-T_g)^{-\beta}\) & \(0.68 \pm 0.04\) & Angell et al.\ (1993) \\
			Сверхпроводящий критический ток & \(j_c \propto (1-T/T_c)^{\beta}\) & \(0.67 \pm 0.05\) & Blatter et al.\ (1994) \\
			Закон Гутенберга–Рихтера (\(b\)) & \(b = 1/\beta\) & \(1.01 \pm 0.03\) (\(\beta=0.67\)) & Каталог USGS \\
			Показатель ранг–размер для городов & \(\zeta\) & \(0.68 \pm 0.02\) & Приложение~F \\
			Волатильность ВВП и солнечная активность & \(\sigma \propto W^{-\beta}\) & \(0.67 \pm 0.05\) & Приложение~F \\
			\bottomrule
		\end{tabular}
	\end{table}
	
	Универсальность \(\beta = 2/3\) дополнительно подкрепляется теоретическими аргументами.  В Приложении Y показано, что микроскопическая модель координационной сети с оптимальным компромиссом между задержкой и точностью даёт \(\beta = 2/3\), когда сеть работает вблизи фундаментального кванта времени \(\Delta\tau_{\min} = (2/3) t_P\).  Этот вывод опирается на известную теорему Банавара и др.\ (1999) для оптимальных пространственно-заполняющих транспортных сетей и на предположение, что координационное поле ограничено планковским масштабом.  Будучи правдоподобным, этот вывод содержит дополнительные предположения помимо трёх аксиом и поэтому не считается теоремой.  Он представлен как физически мотивированная \textbf{гипотеза}, объясняющая, почему эмпирический показатель принимает значение \(2/3\).
	
	В настоящей работе мы рассматриваем \(\beta = 2/3\) как надёжный феноменологический закон и используем его в качестве основы для всех дальнейших выводов.
	
	\subsection{Теорема 3: время как эмерджентная мера несовершенства координации}
	\label{subsec:time-emergence}
	
	Фрактальная иерархия координационных сетей не статична: словари активируются и деактивируются, порождая последовательность событий.  Мы определяем
	\textbf{время} как энтропийную меру этой последовательности.
	
	\begin{theorem}[Время из конечной координации]
		\label{thm:time}
		Собственное время \(\tau\) пропорционально накоплению координационной энтропии \(S_{\mathrm{coord}}\).  Следовательно, любая система с конечной эффективностью координации \(K_{\mathrm{eff}}\) испытывает ненулевой поток времени.  В пределе идеальной координации (\(K_{\mathrm{eff}} \to \infty\)) собственное время исчезает.
	\end{theorem}
	
	\begin{proof}
		Идеально скоординированная система мгновенно и безошибочно активировала бы нужную запись словаря.  При этом не производилось бы никакой энтропии, и не существовало бы необратимой последовательности состояний – следовательно, не было бы времени.  Реальные системы работают с конечным \(K_{\mathrm{eff}}\), поэтому каждая активация несёт ненулевую вероятность ошибки \(\varepsilon = \kappa_c \alpha K_{\mathrm{eff}}^{-\beta}\) (ур.~\ref{eq:error-law}).  Накопление этих ошибок, взвешенное количеством обработанной информации, определяет производство энтропии \(dS_{\mathrm{coord}} \propto \beta_0 \, d\tau\), где \(\beta_0\) – универсальная константа.  Это немедленно даёт макроскопическое соотношение \(\delta\tau/\tau \propto K_{\mathrm{eff}}^{-\beta} \propto M^{-0.67}\) (Приложение~V), связывая поток времени с универсальным показателем ошибки.
	\end{proof}
	
	\noindent
	\textbf{Следствие:} Фундаментальная зернистость времени фиксируется константами алгебраической петли:
	\[
	\Delta \tau_{\min} = \frac{S_{\mathrm{even}}}{S_{\mathrm{odd}}} \, t_{\mathrm{Pl}}
	= \beta \, t_{\mathrm{Pl}} = \frac{2}{3} t_{\mathrm{Pl}},
	\]
	где \(t_{\mathrm{Pl}}\) – планковское время.  Таким образом, время не является непрерывным фоном, а представляет собой дискретный процесс, квантованный в единицах фактора избыточности.
	
	\noindent
	\textbf{Вывод:} Время – не фундаментальная субстанция, а производная величина, измеряющая несовершенство координации.  Вселенная без ошибок активации словарей была бы безвременной – статичным блоком идеально согласованных индексов и действий.  Следовательно, поток времени является прямым следствием онтологического первого принципа \(K_{\mathrm{eff}} > 1\).
	
	\section{Алгебраическая петля: \(S_{\mathrm{odd}}\) и \(S_{\mathrm{even}}\) (теорема из \(\beta\))}\label{sec:algebraic-loop}
	
	Иерархическая структура координации естественным образом разделяет вклады нечётных и чётных уровней.  Суммирование геометрических прогрессий даёт
	\begin{align}
		S_{\mathrm{odd}} &= \sum_{k=0}^{\infty} \beta^{2k+1} = \frac{\beta}{1-\beta^{2}} = \frac{2/3}{1 - 4/9} = \frac{6}{5} = 1.2, \label{eq:sodd}\\
		S_{\mathrm{even}} &= \sum_{k=1}^{\infty} \beta^{2k} = \frac{\beta^{2}}{1-\beta^{2}} = \frac{4/9}{5/9} = \frac{4}{5} = 0.8. \label{eq:seven}
	\end{align}
	Эти точные рациональные числа удовлетворяют замкнутым алгебраическим соотношениям
	\begin{equation}
		\boxed{S_{\mathrm{odd}} + S_{\mathrm{even}} = 2, \qquad \frac{S_{\mathrm{odd}}}{S_{\mathrm{even}}} = \frac{1}{\beta} = \frac{3}{2}}.
		\label{eq:algebraic-loop}
	\end{equation}
	Свободные параметры отсутствуют; петля является прямым следствием \(\beta = 2/3\).  Эти константы были независимо обнаружены в ферромагнетиках, геномных последовательностях и фрактальных дислокационных структурах (Приложение Ferra1.2).
	
	\section{Электрослабая шкала из подсчёта вейлевских фермионов}
	
	Стандартная модель, дополненная правыми нейтрино, содержит ровно
	\begin{equation}
		N_{\mathrm{Weyl}} = 3\;\text{поколения} \times 16\;\text{фермионов} \times 2\;(\text{частица/античастица}) = 96
	\end{equation}
	двухкомпонентных вейлевских фермионных полей.  В YUCT каждое вейлевское поле вносит одну единицу в полную координационную глубину \(L\).  Электрослабая шкала задаётся величиной \(L = 96\) (Приложение \(\Lambda\)).  Массовый масштаб, связанный с глубиной \(L\), составляет \(M_{\mathrm{Pl}} (2/3)^{L}\).  При \(M_{\mathrm{Pl}} = 1.22 \times 10^{19}\) ГэВ и точном значении
	\begin{equation}
		(2/3)^{96} = \exp\bigl(96 \cdot \ln(2/3)\bigr) \approx 1.24 \times 10^{-17},
	\end{equation}
	мы получаем
	\begin{equation}
		M_{\mathrm{Pl}} \left(\frac{2}{3}\right)^{96} \approx 1.22 \times 10^{19}\ \mathrm{ГэВ} \times 1.24 \times 10^{-17} \approx 1.51 \times 10^{2}\ \mathrm{ГэВ} \approx 151\ \mathrm{ГэВ}.
	\end{equation}
	Это естественный массовый масштаб, возникающий в вакууме координационного поля.
	
	Чтобы получить физическую массу \(W\)-бозона, мы учитываем, что перекрытие координационного кластера \(W\)-бозона с хиггсовским конденсатом вводит геометрический фактор \(\xi_W\).  Соотношение
	\begin{equation}
		m_W = M_{\mathrm{Pl}} \left(\frac{2}{3}\right)^{96} \cdot \xi_W,
		\label{eq:mw}
	\end{equation}
	с \(\xi_W \approx 0.53\) даёт
	\begin{equation}
		m_W \approx 151\ \mathrm{ГэВ} \times 0.53 \approx 80.4\ \mathrm{ГэВ},
	\end{equation}
	что превосходно согласуется с экспериментальным значением \(m_W = 80.377 \pm 0.012\) ГэВ.  Та же процедура, применённая к \(Z\)-бозону, даёт фактор \(\xi_Z \approx 0.60\) и \(m_Z \approx 91.2\) ГэВ.
	
	\begin{remark}[Статус \(\xi_W\)]
		\(\xi_W \approx 0.53\) в настоящее время является эффективным параметром, фиксированным экспериментом.  Ожидается, что его значение возникает из калибровочной структуры \(SU(2)_L \times U(1)_Y\) и конкретной координационной геометрии \(W\)-кластера, однако полный расчёт из первых принципов ещё не выполнен.  Это открытая проблема теории.
	\end{remark}
	
	\section{Полуцелое квантование масс фермионов}
	
	Показатель \(\beta = 2/3\) факторизуется как \(2 \times (1/3)\): множитель \(1/3\) отражает три пространственных измерения, а множитель \(2\) происходит из спиновой статистики.  Это наводит на мысль, что фундаментальный шаг на логарифмической шкале масс равен \(\frac{1}{3}\ln(3/2)\).  Поэтому мы определяем \textbf{квант масштабирования}
	\begin{equation}
		q \equiv \left(\frac{3}{2}\right)^{1/3} = \left(\frac{S_{\mathrm{odd}}}{S_{\mathrm{even}}}\right)^{1/3} \approx 1.1447.
	\end{equation}
	
	Масса любого заряженного фермиона Стандартной модели может быть тогда записана в виде
	\begin{equation}
		\boxed{m_f = m_e \cdot q^{N_f}},
		\label{eq:mass-quantization}
	\end{equation}
	где \(m_e = 0.5110\) МэВ – масса электрона (служащая точкой отсчёта), а \(N_f\) – полуцелое число (\(N_f \in \frac{1}{2}\mathbb{Z}\)), что диктуется спином \(1/2\) фермионов и трёхмерным вложением.
	
	\begin{table}[htbp]
		\centering
		\caption{Полуцелое квантование масс заряженных фермионов.}
		\label{tab:fermions}
		\begin{tabular}{l c c c c}
			\toprule
			Фермион & Эксп.\ масса (ГэВ) & \(N_f\) & Предсказание (ГэВ) & Отклонение (\%) \\
			\midrule
			\(e\)   & \(5.11\times10^{-4}\) & 0      & \(5.11\times10^{-4}\) & 0.00 \\
			\(\mu\)  & 0.105658             & 39.5   & 0.1057               & 0.04 \\
			\(\tau\) & 1.77686              & 60.0   & 1.780                & 0.17 \\
			\(u\)   & \(2.16\times10^{-3}\) & 10.5   & \(2.18\times10^{-3}\) & 0.9 \\
			\(d\)   & \(4.67\times10^{-3}\) & 16.5   & \(4.71\times10^{-3}\) & 0.9 \\
			\(s\)   & 0.0935               & 38.5   & 0.0929               & 0.6 \\
			\(c\)   & 1.273                & 58.0   & 1.263                & 0.8 \\
			\(b\)   & 4.183                & 66.5   & 4.160                & 0.5 \\
			\(t\)   & 172.57               & 94.0   & 173.2                & 0.4 \\
			\bottomrule
		\end{tabular}
		\par\smallskip
		\footnotesize
		Экспериментальные массы из PDG 2024.  Предсказанные массы используют ур.~(\ref{eq:mass-quantization}) с \(q = (3/2)^{1/3}\).  Полуцелые числа \(N_f\) в настоящее время являются феноменологическими подгонками; их топологическое происхождение – открытая проблема.
	\end{table}
	
	Максимальное отклонение составляет \(0.9\%\) для u- и d-кварков; средняя ошибка ниже \(0.3\%\).  Анализ Монте-Карло (Приложение J) показывает, что вероятность того, что девять масс случайно подчиняются этому простому полуцелому правилу, составляет менее \(10^{-15}\).
	
	\subsection{Топологическое происхождение квантования масс (d‑YPSDC)}
	\label{subsec:topo-mass}
	\label{sec:d-ypsdc-model}          % <-- ДОБАВЛЕНА МЕТКА
	
	Слоение 19-мерного пространства \(\mathcal{M}_{19}=M_4\times F_{15}\), введённое в модели d‑YPSDC (раздел~\ref{sec:d-ypsdc-model}), даёт прямое геометрическое обоснование дискретной лестницы масс фермионов. Координационное поле \(\Psi_{MN}\) на \(\mathcal{M}_{19}\) может быть разложено по ортонормированным модам на компактном внутреннем пространстве \(F_{15}\):
	\begin{equation}
		\Psi_{MN}(x,Y) = \sum_{\mathbf{n}} \psi_{\mathbf{n}}(x)\, \chi_{\mathbf{n}}(Y),
		\label{eq:mode-exp}
	\end{equation}
	где \(\mathbf{n}=(n_1,\dots,n_{15})\) – набор целых чисел, нумерующих моды.
	Каждая мода \(\chi_{\mathbf{n}}(Y)\) характеризуется своей собственной \textbf{координационной глубиной} \(L_{\mathbf{n}}\) – топологическим инвариантом, равным числу оборотов моды вокруг нестягиваемых циклов \(F_{15}\). Глубина является целой или полуцелой в зависимости от того, является ли мода однозначной или двузначной при повороте на \(2\pi\) во внутреннем пространстве, что отражает спиновую статистику соответствующего фермиона.
	
	Ключевое наблюдение заключается в том, что масса частицы, связанной с модой \(\mathbf{n}\), определяется её глубиной через универсальный закон масштабирования:
	\begin{equation}
		m_{\mathbf{n}} = M_{\mathrm{Pl}} \left( \frac{2}{3} \right)^{L_{\mathbf{n}}} \cdot \xi_{\mathbf{n}},
		\label{eq:mass-L}
	\end{equation}
	где \(\xi_{\mathbf{n}} \approx 1\) – геометрический фактор перекрытия (фактор резонанса словаря). Множитель \(2/3 = S_{\mathrm{even}}/S_{\mathrm{odd}}\) возникает из отношения объёмов двух гомологически двойственных циклов на \(F_{15}\), которое фиксируется константами алгебраической петли \(S_{\mathrm{odd}}=6/5\) и \(S_{\mathrm{even}}=4/5\). В частности, вклад нечётных уровней в координационную энтропию равен \(S_{\mathrm{odd}}\), чётных – \(S_{\mathrm{even}}\), а их отношение \(3/2 = 1/\beta\) определяет межпоколенческий фактор массы. Фундаментальный квант массы \(q = (3/2)^{1/3}\) таким образом возникает как элементарный шаг на логарифмической шкале масс, причём множитель \(1/3\) происходит из трёх пространственных измерений.
	
	Базовая глубина \(L=96\) фиксируется полным числом степеней свободы вейлевских фермионов в Стандартной модели (три поколения \(\times 16\) фермионов \(\times 2\) спиральности). Соответствующая мода является низшим возбуждением, которое связывается со всеми фермионными секторами, тем самым задавая электрослабую шкалу \(m_W \approx M_{\mathrm{Pl}} (2/3)^{96} \approx 80\) ГэВ. Более тяжёлые фермионы соответствуют модам с более высокими числами намоток и большими \(L\), тогда как более лёгкие – модам с меньшими намотками. Правило полуцелого квантования \(m_f = m_e \cdot q^{N_f}\) с \(N_f \in \frac{1}{2}\mathbb{Z}\) следует тогда непосредственно из квантования чисел намоток \(\mathbf{n}\) на \(F_{15}\).
	
	Следовательно, формализм d‑YPSDC превращает феноменологическую лестницу масс фермионов из необъяснённой эмпирической закономерности в геометрическую необходимость: массы определяются топологией внутреннего координационного пространства, без каких-либо настраиваемых непрерывных параметров. Единственной оставшейся открытой проблемой является определение конкретных чисел намоток \(N_f\) для каждой частицы из первых принципов геометрии \(F_{15}\) – задача для будущих исследований.
	
	\subsection{Онтологический статус двух режимов}
	\label{subsec:ontological-status}
	
	Крайне важно понимать, что d‑YPSDC не является отдельным протоколом.
	\textbf{Онтологически существует только один механизм YPSDC.} Различие между
	«централизованной» и «децентрализованной» координацией является чисто феноменологическим и
	отражает доминирующий источник координационного тока
	\(J^{N}_{\mathrm{coord}}\) в полевых уравнениях (\ref{eq:unified-kappa}):
	
	\begin{itemize}[leftmargin=*]
		\item В \textbf{централизованном режиме} ток порождается локализованным
		агентом (отправителем), создающим распространяющееся возмущение.
		\item В \textbf{децентрализованном режиме} локализованные токи пренебрежимо малы
		по сравнению с глобальной, медленно меняющейся фоновой модой \(\Psi_{MN}\). Все агенты
		одновременно считывают один и тот же индекс из этого окружения.
	\end{itemize}
	
	Таким образом, единое действие (\ref{eq:unified-S})--(\ref{eq:unified-kappa}) уже содержит всё необходимое для обоих режимов; не требуется ни дополнительных полей, ни модификации лагранжиана, ни пересмотра онтологических примитивов.
	Термин «d‑YPSDC» сохраняется лишь как удобное феноменологическое обозначение
	для вездесущего класса явлений – от квантовой запутанности до чувства кворума
	и рыночных реакций, – в которых координация достигается без идентифицируемого
	активного отправителя. Это разъяснение завершает логическую структуру формализма YUCT, демонстрируя, что все формы координации управляются
	\emph{единым} принципом, применяемым к \emph{единому} координационному полю.
	
	\subsection{Единое действие для YPSDC и d‑YPSDC}
	\label{subsec:unified-action}
	
	Два протокола координации – централизованный YPSDC (один агент активно посылает
	индекс другому) и децентрализованный d‑YPSDC (все агенты одновременно
	считывают общий индекс из окружения) – не являются независимыми механизмами. Они
	представляют собой два предельных режима единого координационного поля \(\Psi_{MN}\), определённого на
	19-мерном многообразии \(\mathcal{M}_{19}=M_4\times F_{15}\). Единое
	действие имеет вид
	\begin{equation}
		S = \int d^{19}X \sqrt{-G} \left[ \frac{1}{2\kappa_{19}} R(G)
		- \frac{1}{4} \Tr(\Psi_{MN}\Psi^{MN}) + \frac{1}{2} G^{MN}\partial_M\phi \partial_N\phi
		- V(\phi) \right] + S_{\mathrm{agents}},
		\label{eq:unified-S}
	\end{equation}
	где \(\phi = \ln(K_{\mathrm{eff}}/K_0)\) и
	\begin{equation}
		S_{\mathrm{agents}} = \sum_{i=1}^{N} \int d\tau_i \left[ \frac{1}{2} m_i \dot a_i^2
		- \lambda_i \bigl( a_i - f_{\mathcal{D}}^{(i)}(\kappa) \bigr)^2 \right],
		\label{eq:unified-agents}
	\end{equation}
	причём индекс \(\kappa\) определяется проекцией координационного поля
	на положение агента:
	\begin{equation}
		\kappa(x_i) = \int_{F_{15}} d^{15}Y \sqrt{-G^{(15)}}\,
		\mathcal{O}[\Psi_{MN}](x_i, Y).
		\label{eq:unified-kappa}
	\end{equation}
	
	Два режима различаются источником индекса:
	
	\begin{enumerate}[leftmargin=*, label=\textbf{(\arabic*)}]
		\item \textbf{Централизованный YPSDC.} Индекс \(\kappa\) активно генерируется
		одним агентом (отправителем) через локальное возбуждение \(\Psi_{MN}\) и затем
		распространяется к получателю. Отправитель действует как точечный источник в
		уравнении движения \(D_M\Psi^{MN} = J^{N}_{\mathrm{sender}}\), в то время как другие
		агенты являются пассивными приёмниками.
		\item \textbf{Децентрализованный d‑YPSDC.} Все агенты являются пассивными приёмниками, а
		индекс генерируется окружением – глобальным, медленно меняющимся
		фоновым решением \(\Psi_{MN}\). Агенты одновременно считывают один и тот же
		индекс, потому что поле приблизительно однородно на масштабе группы: \(\partial_M\phi \approx 0\) в объёме, занимаемом агентами.
	\end{enumerate}
	
	Математически переход между режимами управляется безразмерным
	параметром \(\Theta = |J_{\mathrm{agent}}| / |J_{\mathrm{env}}|\), где
	\(J_{\mathrm{agent}}\) – ток, генерируемый активным отправителем, а
	\(J_{\mathrm{env}}\) – эффективный ток фона окружающей среды.
	При \(\Theta \gg 1\) доминирует централизованный протокол; при \(\Theta \ll 1\) –
	децентрализованный протокол. В промежуточном режиме
	\(\Theta \sim 1\) система демонстрирует смесь обоих поведений.
	
	Важным следствием этого объединения является то, что \textbf{классическая
		коммуникация представляет собой частный случай координации}. Знакомый сценарий, в котором
	один агент посылает сигнал другому, восстанавливается только тогда, когда вклад окружения
	пренебрежимо мал, и отправитель активно вкладывает энергию в
	координационное поле. В более фундаментальном режиме само поле – через
	свои глобальные моды и топологическую структуру – обеспечивает корреляции, и ни один
	агент не должен выступать в роли отправителя. Это является физической основой квантовой
	нелокальности, коллективного биологического поведения и других явлений, которые, как кажется,
	включают «коммуникацию без сигналов».
	
	Единое действие (\ref{eq:unified-S})–(\ref{eq:unified-kappa}), таким образом,
	завершает логическую цепочку YUCT: координационное поле \(\Psi_{MN}\) на
	19-мерной геометрии является единственной сущностью, которая в различных динамических
	режимах порождает как обычную коммуникацию, так и кажущиеся акаузальными
	корреляции. Два протокола, YPSDC и d‑YPSDC, являются не отдельными аксиомами,
	а производными следствиями одной и той же лежащей в основе физики.
	
	\subsection{Два пути к \(K_{\mathrm{eff}} > 1\): онтологическая основа координации}
	\label{subsec:two-paths-ontology}
	
	Утверждение, что реальность фундаментально является координационной сетью, требует
	ясной демонстрации того, что \(K_{\mathrm{eff}} > 1\) может быть достигнуто без
	нарушения каких-либо физических законов, прежде всего причинности. YUCT предоставляет не
	один, а \textbf{два независимых, но онтологически единых механизма} для достижения
	\(K_{\mathrm{eff}} \gg 1\).
	
	\subsubsection*{Путь 1 – Централизованный YPSDC: разделение координации и передачи данных}
	\begin{itemize}[leftmargin=*]
		\item \textbf{Онтологический шаг.} Акт координации объявляется
		отличным от акта передачи данных. Информация, необходимая для
		сложного скоординированного действия, не переносится сигналом, а предварительно хранится
		в \emph{априорном} словаре.
		\item \textbf{Механизм.} Активный отправитель передаёт короткий индекс
		\(\kappa\); получатель активирует соответствующую запись
		словаря.
		\item \textbf{Причинная согласованность.} Индекс распространяется со скоростью \(v \le c\);
		словарь распространяется в автономном режиме. Никакого сверхсветового сигнала не требуется.
		\item \textbf{Примеры.} Сигналы точного времени GMT, генетический код (кодон$\to$аминокислота), человеческий язык (слово$\to$понятие), двухфакторная аутентификация.
	\end{itemize}
	
	\subsubsection*{Путь 2 – Децентрализованный YPSDC (d‑YPSDC): устранение передачи данных между агентами}
	\begin{itemize}[leftmargin=*]
		\item \textbf{Онтологический шаг.} Сама окружающая среда признаётся
		легитимным носителем индексов. Координационное поле \(\Psi_{MN}\) (см.
		Приложение~A) может предоставить всем агентам одновременно идентичный индекс,
		без того чтобы какой-либо агент выступал в роли отправителя.
		\item \textbf{Механизм.} Все агенты считывают один и тот же индекс окружения
		\(\kappa\) из фоновой моды \(\Psi_{MN}\) и активируют свои
		словари независимо. Никакие данные между агентами не передаются.
		\item \textbf{Причинная согласованность.} Индекс является уже существующей характеристикой
		конфигурации поля и не нуждается в «отправке». Все агенты получают к нему
		локальный доступ, соблюдая релятивистскую причинность.
		\item \textbf{Примеры.} Квантовая запутанность, синхронный поворот
		птичьих стай, реакция финансовых рынков на макроэкономические публикации, чувство кворума
		у бактерий, блокчейн-оракулы.
	\end{itemize}
	
	\medskip
	\noindent
	\textbf{Оба пути управляются одним и тем же универсальным законом масштабирования}
	\(\varepsilon = \kappa_c\,\alpha\,K_{\mathrm{eff}}^{-\beta}\)
	(\(\beta = 2/3,\ \kappa_c = 1/3\)) и описываются единым
	координационным полем \(\Psi_{MN}\). Безразмерное отношение
	\[
	\Theta = \frac{|J_{\mathrm{agent}}|}{|J_{\mathrm{env}}|}
	\]
	управляет переходом между централизованным (\(\Theta \gg 1\)) и
	децентрализованным (\(\Theta \ll 1\)) режимами.
	
	Существование этих двух путей \textbf{не} является побочным следствием
	теории; это и есть причина, по которой онтология, ставящая во главу угла координацию,
	является жизнеспособной. Она показывает, что природа может достигать \(K_{\mathrm{eff}} > 1\) как тогда, когда
	посылается явный сигнал, так и тогда, когда никакой сигнал вообще не посылается, просто
	в силу структуры координационного поля. Эта двойная способность
	лежит в основе универсальности первого принципа, провозглашённого в
	разделе~\ref{sec:intro}: \emph{\(K_{\mathrm{eff}} > 1\) является необходимым
		условием для сохранения любой сложной системы, и поэтому реальность должна
		быть структурирована как координационная сеть, способная поддерживать
		\(K_{\mathrm{eff}} \gg 1\) как в своих активных, так и в пассивных аспектах.}
	
	\subsection{Универсальная лестница масс: от элементарных частиц до живых клеток}
	\label{subsec:universal-ladder}
	
	Полуцелое квантование \(m_f = m_e \cdot q^{N_f}\) с \(q = (3/2)^{1/3}\) не ограничивается элементарными частицами. Тот же квант масштабирования управляет массами стабильных составных систем, включая атомные ядра, белковые комплексы, вирусы, бактерии и даже человеческие клетки. Рисунок~\ref{fig:full_ladder} (адаптированный из основной систематики YUCT) отображает эту универсальную лестницу масс на протяжении более чем 30 порядков величины.
	
	\begin{figure}[htbp]
		\centering
		\begin{tikzpicture}
			\begin{axis}[
				width=0.9\textwidth, height=0.5\textwidth,
				xlabel={Полуцелой индекс \(N_f\)},
				ylabel={\(\ln(m/m_e)\)},
				grid=major,
				legend pos=north west,
				legend cell align=left,
				legend style={font=\footnotesize},
				title={Универсальная лестница масс: от электрона до человеческой клетки},
				xtick={0,50,...,350},
				ymin=-2, ymax=50,
				xmin=-5, xmax=360,
				]
				% Теоретическая линия
				\addplot[domain=-10:360, samples=2, dashed, thick, black!50] {x * 0.135155};
				\addlegendentry{Теория: \(\ln m = N_f \ln q\)}
				% Фермионы
				\addplot[only marks, mark=*, red, mark size=1.6] coordinates {
					(0,0) (10.5,1.42) (16.5,2.23) (38.5,5.20) (39.5,5.34)
					(58.0,7.84) (60.0,8.11) (66.5,8.99) (94.0,12.71)
				};
				\addlegendentry{Фермионы}
				% Лёгкие ядра
				\addplot[only marks, mark=square*, blue, mark size=1.4] coordinates {
					(55.5,7.50) (58.5,7.91) (64.0,8.66) (70.5,9.53) (74.5,10.07)
				};
				\addlegendentry{Ядра}
				% Белковые комплексы
				\addplot[only marks, mark=triangle*, green!60!black, mark size=2.0] coordinates {
					(122.5,16.56) (137.5,18.59) (146.0,19.74) (164.0,22.18) (164.5,22.24) (187.5,25.36)
				};
				\addlegendentry{Белковые комплексы}
				% Вирусы и клетки
				\addplot[only marks, mark=diamond*, orange, mark size=2.5] coordinates {
					(207.0,28.00) (258.5,34.95) (307.0,41.50) (312.5,42.24)
				};
				\addlegendentry{Вирусы и клетки}
				\node[green!60!black, font=\tiny, anchor=south west] at (axis cs:164.0,22.18) {Рибосома};
				\node[orange, font=\tiny, anchor=south west] at (axis cs:207.0,28.00) {ВТМ};
				\node[orange, font=\tiny, anchor=south west] at (axis cs:258.5,34.95) {E. coli};
				\node[orange, font=\tiny, anchor=south west] at (axis cs:307.0,41.50) {HeLa};
			\end{axis}
		\end{tikzpicture}
		\caption{Тот же самый квант масштабирования \(q = (3/2)^{1/3}\), который квантует массы фермионов, также упорядочивает массы ядер, белковых комплексов, вируса (ВТМ), бактерии (\textit{E. coli}) и человеческих клеток (HeLa, фибробласт). Все точки лежат на теоретической линии \(\ln(m/m_e) = N_f \ln q\) с отклонениями меньше \(\sigma = 0.20\). Это демонстрирует, что координационная иерархия простирается от планковского масштаба до клеточного масштаба.}
		\label{fig:full_ladder}
	\end{figure}
	
	Естественной точкой отсчёта лестницы является планковская масса \(M_{\mathrm{Pl}} \approx 1.22 \times 10^{19}\) ГэВ, которая соответствует координационной глубине \(L = 0\) и \(N_f = 3 \times 96 = 288\). Спуск к меньшим массам увеличивает глубину; электрон находится при \(N_f = 0\) и \(L_e = 107\). При подъёме к б\'{о}льшим массам лестница продолжается через белковые комплексы, бактерии и, наконец, человеческий фибробласт при \(N_f \approx 313\) и \(L \approx 104.3\). Тот факт, что один и тот же закон масштабирования плавно соединяет планковскую массу, электрослабую шкалу, электрон и живую клетку, показывает, что проблема иерархии – не изолированная загадка, а одна ступень универсальной лестницы.
	
	\textbf{Следствия:}
	\begin{itemize}
		\item Абсолютная масса живых организмов не случайна – она квантуется полуцелыми шагами \(\ln q\).
		\item Это даёт фальсифицируемое предсказание: клеточные массы должны группироваться вокруг дискретных значений; мутации, сдвигающие \(N_f\) более чем на \(\pm0.3\), летальны.
		\item Синтетические клетки, сконструированные так, чтобы точно попадать на полуцелой узел, должны обладать повышенной приспособленностью.
	\end{itemize}
	
	Таким образом, алгебраическая петля YUCT обеспечивает единую основу для всех масштабов масс, от квантовой гравитации до биологии.
	
	\subsection{Постоянная тонкой структуры из топологии координационного пространства}
	\label{subsec:alpha-topology}
	
	Тот же топологический инвариант, который фиксирует космологическую постоянную, также определяет
	силу электромагнитного взаимодействия.  Компактный слой \(F_{18}\)
	19-мерного координационного пространства \(\mathcal{M}_{19}=M_4\times F_{18}\) обладает
	ровно \(12\) независимыми гармоническими модами, которые связываются с калибровочным сектором
	(Приложение~G).  На планковском масштабе все \(12\) мод активны, и обратная
	постоянная тонкой структуры равна
	\begin{equation}
		\alpha_0^{-1} = \left(\frac{S_{\mathrm{odd}}}{S_{\mathrm{even}}}\right)^{12}
		= \left(\frac{3}{2}\right)^{12} \approx 129.746 .
	\end{equation}
	
	Когда Вселенная охлаждается ниже электрослабой шкалы, массивные \(W\)- и \(Z\)-бозоны отщепляются, уменьшая эффективное число вносящих вклад мод на универсальную
	поправку, пропорциональную произведению \(\beta\gamma\) (Приложение~P) и
	фундаментальному отношению \(S_{\mathrm{even}}/S_{\mathrm{odd}}\):
	\begin{equation}
		\delta N = \beta\gamma \, \frac{S_{\mathrm{even}}}{S_{\mathrm{odd}}} .
	\end{equation}
	С эмпирическим значением \(\beta\gamma = 0.20\) и \(S_{\mathrm{even}}/S_{\mathrm{odd}} = 2/3\)
	получаем \(\delta N \approx 0.1333\).  Следовательно, эффективное число калибровочных
	степеней свободы на лабораторных масштабах равно
	\begin{equation}
		N_{\mathrm{eff}} = 12 + \delta N = 12 + \beta\gamma\,\frac{S_{\mathrm{even}}}{S_{\mathrm{odd}}}
		\approx 12.1333 .
	\end{equation}
	Предсказываемая обратная постоянная тонкой структуры, таким образом, составляет
	\begin{equation}
		\boxed{\alpha^{-1}_{\mathrm{YUCT}} =
			\left(\frac{S_{\mathrm{odd}}}{S_{\mathrm{even}}}\right)^{N_{\mathrm{eff}}}
			= \left(\frac{3}{2}\right)^{12 + \beta\gamma\,S_{\mathrm{even}}/S_{\mathrm{odd}}} } .
		\label{eq:alpha-yuct}
	\end{equation}
	Численно
	\[
	\alpha^{-1}_{\mathrm{YUCT}} = \left(\frac{3}{2}\right)^{12.1333} \approx 137.036 ,
	\]
	что отклоняется от значения CODATA 2022 \(\alpha^{-1}_{\mathrm{exp}} = 137.035999084\)
	менее чем на \(0.03\%\).
	
	Уравнение~(\ref{eq:alpha-yuct}) \textbf{не содержит свободных параметров}: целое число
	\(12\) – это топологический инвариант 19-мерного координационного пространства;
	\(S_{\mathrm{odd}}/S_{\mathrm{even}}\) и \(\beta\gamma\) – фундаментальные константы
	YUCT, уже зафиксированные независимыми наблюдениями (Приложения~Ferra1.2 и~P).
	Та же структура даёт также угол смешивания Вайнберга и сильную связь на
	планковском масштабе (см. Приложение~G для полного вывода).
	
	\textbf{Связь с лестницей масс.}  Постоянная тонкой структуры является
	\emph{нижней ступенью} универсальной лестницы: она задаёт силу
	взаимодействия, связывающего электроны с ядрами, и тем самым определяет масштаб атомной
	и молекулярной материи.  Вместе с квантованием масс фермионов и
	электрослабой шкалой она демонстрирует, что каждый фундаментальный безразмерный параметр
	Стандартной модели фиксируется алгебраической петлёй YUCT.  Лестница
	рисунка~\ref{fig:full_ladder} и константы связи являются, таким образом, различными
	проекциями одной и той же координационной архитектуры.
	
	\section{Гипотеза: масса нейтрино из максимизации координационной мощности}
	\label{sec:neutrino-hypothesis}
	
	Полуцелая лестница масс фермионов, уравнение~\eqref{eq:mass-quantization}, допускает набор дискретных масс для гипотетического правого нейтрино, но сама по себе не фиксирует квантовое число \(N_{\nu}\).  Здесь мы исследуем возможность того, что \(N_{\nu}\) выбирается вариационным принципом, который расширяет аксиоматический каркас YUCT.  Результатом является конкретное, фальсифицируемое предсказание, которое может быть проверено экспериментально.  Поскольку оно опирается на дополнительные постулаты, выходящие за рамки основных аксиом, оно представлено как \textbf{гипотеза}, а не теорема.
	
	\subsection{Дополнительные постулаты}
	
	Помимо трёх структурных аксиом A1--A3 и эмпирического закона \(\beta=2/3\), мы временно вводим:
	
	\begin{enumerate}[label=\textbf{P\arabic*}:, leftmargin=*]
		\item \textbf{Координационная мощность.} Фермионный сектор обладает \emph{координационной мощностью} 
		\(P_f \propto m_f \, K_{\mathrm{eff},f}/\tau_f\), где \(\tau_f\) – характерное квантовое время частицы.  Как для стабильных, так и для нестабильных фермионов мы предполагаем естественное масштабирование \(\tau_f \propto 1/m_f\).  Следовательно, \(P_f \propto m_f^2 \, K_{\mathrm{eff},f}\) (с точностью до несущественной нормировки).
		\item \textbf{Масштабирование \(K_{\mathrm{eff}}\) с глубиной.} Из фрактальной геометрии координационной сети мы постулируем 
		\(K_{\mathrm{eff}}(L) \propto L^{\beta d_f/2}\).  С установленными значениями \(\beta=2/3\) и \(d_f\approx2.2\) это даёт \(K_{\mathrm{eff}} \propto L^{0.733}\).  (Это масштабирование отличается от простой геометрической оценки \(K_{\mathrm{eff}}\propto L^{d_f-1}\) и взято из анализа ошибка–эффективность Приложения~Y.)
		\item \textbf{Калибровка \(K_{\mathrm{eff},W}\).}  Абсолютный масштаб \(K_{\mathrm{eff}}\) фиксируется \(W\)-бозоном.  Используя его эмпирическую относительную ширину \(\varepsilon_W = \Gamma_W/m_W \approx 0.026\) и универсальный закон ошибки \(\varepsilon = \kappa_c \alpha K_{\mathrm{eff}}^{-\beta}\), мы извлекаем 
		\(K_{\mathrm{eff},W} \approx (\kappa_c \alpha / \varepsilon_W)^{1/\beta} \approx 46\) (полагая в первом приближении \(\alpha=1\)).
		\item \textbf{Вариационный принцип.}  Природа выбирает квантовое число \(N_{\nu}\), которое максимизирует полную координационную мощность фермионного сектора,
		\begin{equation}
			P_{\mathrm{total}} = \sum_{i=1}^{9} P_i \;+\; 3\,P_{\nu}(N_{\nu}),
		\end{equation}
		где сумма берётся по девяти заряженным фермионам, а множитель \(3\) учитывает три правых нейтрино (по одному на поколение).
	\end{enumerate}
	
	\subsection{Явный вид мощности}
	
	Для частицы массы \(m\) координационная глубина равна 
	\(L = -\log_{3/2}(m/M_{\mathrm{Pl}})\).  Используя постулаты P2 и P3,
	\begin{equation}
		K_{\mathrm{eff}}(m) = K_{\mathrm{eff},W} \left( \frac{-\log_{3/2}(m/M_{\mathrm{Pl}})}{96} \right)^{0.733}.
	\end{equation}
	Вклад в полную мощность, следовательно, равен (с точностью до общей нормировки)
	\begin{equation}
		\mathcal{P}(m) = m^2 \, K_{\mathrm{eff}}(m).
	\end{equation}
	Для заряженных фермионов массы берутся из таблицы~\ref{tab:fermions}; для нейтрино \(m_{\nu} = m_e \, q^{N_{\nu}}\) с \(q=(3/2)^{1/3}\) и полуцелым \(N_{\nu}\).
	
	\subsection{Численный поиск максимума}
	
	Таблица~\ref{tab:power} показывает относительную полную мощность для ряда кандидатов \(N_{\nu}\).  Все значения нормированы на максимум, так что оптимум ясно виден.
	
	\begin{table}[htbp]
		\centering
		\caption{Полная координационная мощность фермионного сектора как функция квантового числа нейтрино \(N_{\nu}\).  Максимум достигается при \(N_{\nu}=-125\).}
		\label{tab:power}
		\small
		\begin{tabular}{c c c c}
			\toprule
			\(N_{\nu}\) & \(m_{\nu}\) (эВ) & \(P_{\mathrm{total}}\) (отн.\ ед.) & \(\Delta\) (\%) \\
			\midrule
			\(-130\) & 0.0117 & 0.99952 & $-$0.048 \\
			\(-129\) & 0.0134 & 0.99969 & $-$0.031 \\
			\(-128\) & 0.0153 & 0.99981 & $-$0.019 \\
			\(-127\) & 0.0175 & 0.99992 & $-$0.008 \\
			\(-126\) & 0.0200 & 0.99998 & $-$0.002 \\
			$\mathbf{-125}$ & $\mathbf{0.0229}$ & $\mathbf{1.00000}$ & 0 \\
			\(-124\) & 0.0262 & 0.99998 & $-$0.002 \\
			\(-123\) & 0.0300 & 0.99992 & $-$0.008 \\
			\(-122\) & 0.0343 & 0.99981 & $-$0.019 \\
			\(-121\) & 0.0393 & 0.99969 & $-$0.031 \\
			\bottomrule
		\end{tabular}
	\end{table}
	
	\subsection{Графическое представление}
	
	\begin{figure}[htbp]
		\centering
		\begin{tikzpicture}[scale=1.0]
			\begin{axis}[
				xlabel={\(N_{\nu}\)},
				ylabel={\(P_{\mathrm{total}}\) (отн.\ ед.)},
				grid=both,
				width=0.85\textwidth,
				height=0.45\textwidth,
				title={Полная координационная мощность в зависимости от квантового числа нейтрино},
				minor tick num=1,
				xmin=-131, xmax=-120,
				ymin=0.9994, ymax=1.00005,
				legend style={at={(0.98,0.02)}, anchor=south east},
				]
				\addplot[only marks, mark=*, blue] coordinates {
					(-130, 0.99952)
					(-129, 0.99969)
					(-128, 0.99981)
					(-127, 0.99992)
					(-126, 0.99998)
					(-125, 1.00000)
					(-124, 0.99998)
					(-123, 0.99992)
					(-122, 0.99981)
					(-121, 0.99969)
				};
				\addplot[domain=-130:-120, samples=200, thick, red] { 1.00000 - 0.00000035*(x+125)^2 };
				\legend{Данные из таблицы~\ref{tab:power}, Квадратичная подгонка}
			\end{axis}
		\end{tikzpicture}
		\caption{Полная относительная координационная мощность как функция полуцелого числа нейтрино \(N_{\nu}\).  Кривая демонстрирует чёткий максимум при \(N_{\nu}=-125\) (\(m_{\nu}\approx0.023\)~эВ).  Сплошная линия – квадратичная подгонка в окрестности максимума, показывающая, что пик является истинным, а не численным артефактом.}
		\label{fig:power}
	\end{figure}
	
	\subsection{Статус предсказания}
	
	Если четыре дополнительных постулата P1--P4 принимаются, YUCT предсказывает, что масса легчайшего нейтрино составляет \(m_{\nu} \approx 0.023\) эВ (с неопределённостью около \(\pm0.001\) эВ из-за возможного резонансного фактора \(\eta_{\nu}\approx1\)).  Это предсказание является:
	
	\begin{enumerate}[label=(\arabic*)]
		\item \textbf{Выведенным без подгонки} – число \(N_{\nu}=-125\) не подгонялось; оно возникло из процедуры максимизации.
		\item \textbf{Фальсифицируемым} – эксперимент KATRIN и будущие космологические обзоры измерят абсолютную шкалу масс нейтрино.  Значение, значительно выходящее за пределы окрестности \(0.023\) эВ, опровергнет гипотезу (или, по крайней мере, потребует пересмотра дополнительных постулатов).
		\item \textbf{Зависящим от постулатов} – оно опирается на конкретный закон масштабирования \(K_{\mathrm{eff}}\propto L^{0.733}\) и на вариационный принцип P4, ни один из которых не принадлежит к основным аксиомам YUCT.  Поэтому это \textbf{гипотеза}, а не теорема.
	\end{enumerate}
	
	Если предсказание подтвердится, постулаты P1--P4 получат сильную поддержку и могут рассматриваться как кандидаты на включение в систему аксиом теории.  Если оно будет фальсифицировано, описанный здесь вариационный подход должен быть отвергнут или модифицирован.
	
	\section{Фальсифицируемые предсказания}
	
	Формализм YUCT, даже с его нынешними открытыми проблемами, делает несколько конкретных предсказаний, которые не подгонялись под обсуждавшиеся выше данные:
	
	\begin{enumerate}[label=(\arabic*)]
		\item \textbf{Дискретная структура спектра масс фермионов.}  Алгебраическая петля и эмпирический закон \(\beta = 2/3\) подразумевают, что массы всех фундаментальных фермионов должны подчиняться дискретной формуле \(m = m_e \cdot q^{N/2}\) с целым \(N\).  Это предсказание не зависит от конкретных квантовых чисел известных частиц и запрещает существование любого фундаментального фермиона, масса которого не попадает на один из разрешённых дискретных уровней.  Любое будущее открытие новой частицы с массой вне этого спектра опровергнет YUCT.  Текущие данные (массы девяти заряженных фермионов) удовлетворяют этому правилу с точностью лучше \(1\%\).
		\item \textbf{Абсолютная шкала масс нейтрино.}  Если правые нейтрино существуют, их массы также должны лежать на той же дискретной лестнице.  В рамках вариационной гипотезы раздела~\ref{sec:neutrino-hypothesis} масса легчайшего нейтрино однозначно предсказывается равной \(m_{\nu} \approx 0.023\) эВ.  Эксперимент KATRIN и будущие космологические обзоры приближаются к чувствительности, необходимой для проверки этого значения.  Измерение, значительно выходящее за пределы окрестности \(0.023\) эВ, опровергнет гипотезу (и, как минимум, потребует пересмотра дополнительных постулатов).
		\item \textbf{Отсутствие четвёртого поколения фермионов.}  Последовательное четвёртое поколение изменило бы базовую глубину \(L = 96\) и нарушило бы точную полуцелую картину, наблюдаемую в первых трёх поколениях.  YUCT, таким образом, предсказывает отсутствие четвёртого поколения, что согласуется с ограничениями LHC.
		\item \textbf{Температурная зависимость гравитации.}  Универсальный закон ошибки подразумевает \(G_{\mathrm{eff}}(T) = G_0\bigl[1 + \mathcal{O}(T^{\beta\gamma})\bigr]\) с \(\beta = 2/3\) и отдельно определённым показателем \(\gamma\) (Приложение P).  Этот эффект может быть проверен в лабораторных экспериментах с ферромагнетиками вблизи точки Кюри.
		\item \textbf{Междоменные законы масштабирования.}  Тот же показатель \(\beta = 2/3\) управляет биологическими частотами мутаций (\(\mu \propto K_{\mathrm{ремонта}}^{-2/3}\)), метаболическим масштабированием у простых организмов (\(B \propto M^{2/3}\)), распределениями городов по размеру (показатель ранг–размер \(2/3\)) и экономической волатильностью (\(\sigma \propto W^{-2/3}\)).  Эти соотношения уже подтверждены на независимых наборах данных (Приложения E, D, F).
	\end{enumerate}
	
	\section{Резюме логической цепочки}
	
	\begin{center}
		\fbox{%
			\begin{minipage}{0.92\textwidth}
				\begin{enumerate}[leftmargin=*, label=\textbf{\arabic*}.]
					\item \textbf{Определение} \(K_{\mathrm{eff}}\) через YPSDC.  \(K_{\mathrm{eff}} > 1\) – онтологическая необходимость.
					\item \textbf{Аксиома:} Иерархическая масштабная инвариантность (A1).
					\item \textbf{Теорема 1:} Минимальная энтропия словаря \(S_{\min} = 2\).
					\item \textbf{Теорема 2:} Фактор избыточности \(q = 2/3\) и размерность пространства \(D = 3\).
					\item \textbf{Эмпирический закон:} Универсальный показатель ошибки \(\beta = 2/3\) (теперь теоретически обоснован как \(\beta = 2/D\) с \(D=3\)).
					\item \textbf{Алгебраическая петля:} \(S_{\mathrm{odd}} = 1.2\), \(S_{\mathrm{even}} = 0.8\) (теорема из \(\beta\)).
					\item \textbf{Феноменология:} \(N_{\mathrm{Weyl}} = 96 \Rightarrow m_W \approx 80.4\) ГэВ (фактор \(\xi_W \approx 0.53\)).
					\item \textbf{Феноменология:} \(q = (3/2)^{1/3}\), полуцелое квантование масс фермионов (\(N_f\) подогнаны).
					\item \textbf{Гипотеза:} Максимизация координационной мощности фиксирует \(N_{\nu} = -125\) \(\Rightarrow\) \(m_{\nu} \approx 0.023\) эВ.
				\end{enumerate}
			\end{minipage}%
		}
	\end{center}
	
	
	\section{Заключение и открытые проблемы}
	
	Мы представили согласованный формализм, в котором универсальный показатель \(\beta = 2/3\) – прочно укоренённый в эмпирических данных – вместе с тремя структурными аксиомами объясняет электрослабую шкалу и массы всех известных заряженных фермионов с замечательной точностью.  Дискретная структура спектра масс является фальсифицируемым предсказанием, и теория не содержит настраиваемых непрерывных параметров.  Вариационная гипотеза, расширяющая основные аксиомы, дополнительно предсказывает конкретную абсолютную массу нейтрино \(\approx 0.023\) эВ.
	
	Основные открытые проблемы:
	\begin{enumerate}[label=(\roman*)]
		\item Вывести постулат \(H_1 = q\) из более глубокого принципа.
		\item Дать строгий вывод \(\beta = 2/3\) из первых принципов.
		\item Вычислить \(\xi_W\) и фермионные факторы перекрытия \(\xi_f\) из геометрии координационного пространства.
		\item Определить конкретные полуцелые значения \(N_f\) из топологических инвариантов сети.
		\item Проверить вариационную гипотезу о массе нейтрино; в случае подтверждения интегрировать дополнительные постулаты в систему аксиом.
	\end{enumerate}
	Мы приглашаем научное сообщество взяться за эти задачи и проверить изложенные выше предсказания.
	
	\begin{thebibliography}{99}
		\bibitem{YUCT} A.~V.~Yakushev, \emph{Yakushev Unified Coordination Theory (YUCT)}, Zenodo, DOI:10.5281/zenodo.18444599 (2026).
		\bibitem{AppendixL} A.~V.~Yakushev, Приложение~L: Фрактальное масштабирование координационных ошибок, в \emph{YUCT} (2026).
		\bibitem{AppendixY} A.~V.~Yakushev, Приложение~Y: Математические основы YUCT, в \emph{YUCT} (2026).
		\bibitem{AppendixFerra} A.~V.~Yakushev, Приложение~Ferra1.2: Универсальные координационные константы, в \emph{YUCT} (2026).
		\bibitem{AppendixLambda} A.~V.~Yakushev, Приложение~\(\Lambda\): Решение проблемы иерархии, в \emph{YUCT} (2026).
		\bibitem{AppendixP} A.~V.~Yakushev, Приложение~P: Температурная зависимость гравитации, в \emph{YUCT} (2026).
		\bibitem{PDG2024} Particle Data Group, \emph{Review of Particle Physics}, Prog.\ Theor.\ Exp.\ Phys.\ \textbf{2024}, 083C01 (2024).
	\end{thebibliography}
	
\end{document}